% !TeX program = lualatex
% !TeX encoding = utf8
% !TeX root = ../Thermod.tex

%\setcounter{chapter}{13}

%=========================================================
\Opensolutionfile{answer}[\currfilebase/\currfilebase-Answers]
\Writetofile{answer}{\protect\section*{\nameref*{\currfilebase}}}
\chapter{Статистичні розподіли. Статистичний зміст ентропії}\label{\currfilebase}
\graphicspath{{\currfilebase/pictures}}
\makeatletter
\def\input@path{{\currfilebase/}}
\makeatother
%=========================================================

%% --------------------------------------------------------
\section{Поняття макро та мікростану системи}
%% --------------------------------------------------------

Тут ми продовжимо тему статистичних розподілів та спробуємо з’ясувати чи можна користуючись статистичним підходом пов’язати макростан системи (тобто значення низки макроскопічних змінних, що характеризують систему) із її мікростаном (розподілом імпульсів, енергій і положень частинок системи) (Рис.~\ref{pic:Pic14.1}).

%---------------------------------------------------------
\begin{figure}[h!]\centering
	\includegraphics[width=0.75\linewidth]{Pic1}
	\caption{Кожному макростану системи, що характеризується низкою макроскопічних параметрів, можуть відповідати різні мікростани}
	\label{pic:Pic14.1}
\end{figure}
%---------------------------------------------------------

Найбільш інтригуючим є питання, чи можна із допомогою статистичного підходу визначити поняття ентропії. З вичерпуючою повнотою ця тема не розкривається навіть у рамках статистичної фізики, тому більшість наших міркувань будуть носити наближений, ілюстративний характер. Заздалегідь зазначимо, що найбільших успіхів статистичний підхід виявляє в системах з слабкою взаємодією між складовими частинами системи. В цьому випадку така характеристика як енергія є адитивною величиною --- при розділенні системи на частини сума енергій складових частин дорівнює енергії цілої системи. Поведінка такого типу систем добре ілюструється моделлю квантових осциляторів, ідеального газу, моделлю спінової системи та інші (Рис.~\ref{pic:Pic14.2}). Серед них, модель ідеального газу має явну перевагу в тому, що для її опису залучаються тільки уявлення класичної механіки Ньютона. Тому надалі усі наші міркування будуть опиратися на цю модель, маючи на увазі, що метод комірок Больцмана застосовуваний до моделі має більш загальне застосування.

%---------------------------------------------------------
\begin{figure}[h!]\centering
	\begin{subfigure}[b]{0.3\textwidth}
		\centering
		\includegraphics[width=\linewidth]{Pic2a}
		\caption{}
	\end{subfigure}
	\hfill
	\begin{subfigure}[b]{0.3\textwidth}
		\centering
		\includegraphics[width=\linewidth]{Pic2b}
		\caption{}
	\end{subfigure}
	\hfill
	\begin{subfigure}[b]{0.3\textwidth}
		\centering
		\includegraphics[width=\linewidth]{Pic2c}
		\caption{}
	\end{subfigure}
	\caption{Приклади систем із слабкою взаємодією: а) спінова система; б) модель квантових осциляторів; в) модель елементарного газу}
	\label{pic:Pic14.2}
\end{figure}
%---------------------------------------------------------

%% --------------------------------------------------------
\section{Фазовий простір координати-імпульс}
%% --------------------------------------------------------

Динамічне описання системи передбачає задання в певний момент часу координат та імпульсів усіх молекул з наступним інтегруванням рівнянь руху. Як відомо, настільки детальний опис руху для термодинамічного опису є надлишковим. У більшості випадків достатньо знати тільки деякі середні характеристики типу $\langle v \rangle$, $v_{\text{ср.кв}}$ та інші так чи інакше зв’язані з макроскопічними параметрами. Приклад з розподілом Максвела показує, що вичерпний опис дає знання розподілу молекул по енергіям.

В загальному випадку, енергія молекул складається з суми кінетичної та потенціальної енергій і в тому шукані розподіли суті функції як швидкостей (імпульсів) так і координат молекул. Таким чином, якщо при виводі розподілу Максвела ми задавалися тільки швидкістю молекули вектору $v$ з точністю до $dv_x$, $dv_y$, $dv_z$, тоді в загальному випадку потрібно вказувати і її координату з певною точністю до $dx$, $dy$, $dz$. Тому введемо шестивимірний фазовий простір, кожна точка якого відповідає певній молекулі і описується координатами $x$, $y$, $z$, $p_x$, $p_y$, $p_z$. Миттєвий стан окремої молекули характеризується певною точкою у цьому просторі, названою фазовою, а динамічний стан усіх $N$ молекул – $N$ фазовими точками. Розіб’ємо увесь фазовий простір на фазові комірки – шестивимірні паралелепіпеди із сторонами $dx$, $dy$, $dz$, $dp_x$, $dp_y$, $dp_z$. Будемо вважати розміри фазової комірки з одного боку досить маленькими, щоб прийняти енергію будь-якої молекули з фазовою точкою, що попадає у $i$-ту комірки $\varepsilon_i$, однаковою для всіх таких молекул, з іншого боку достатньо великими, щоб в кожну таку комірку попадала велика кількість фазових точок (Рис.~\ref{pic:Pic14.3}).

%---------------------------------------------------------
\begin{figure}[h!]\centering
	\begin{subfigure}[b]{0.45\textwidth}
		\centering
		\includegraphics[width=\linewidth]{Pic3a}
		\caption{}
	\end{subfigure}
	\hfill
	\begin{subfigure}[b]{0.45\textwidth}
		\centering
		\includegraphics[width=\linewidth]{Pic3b}
		\caption{}
	\end{subfigure}
	\caption{а) чотирьохвимірний переріз шестивимірного фазового простору координат-імпульсів молекул; б) розміри комірки в цьому просторі мають бути водночас достатньо малими і достатньо великими як у комірки 3, достатніми щоб туди вміщалося багато частинок (виділені червоним)}
	\label{pic:Pic14.3}
\end{figure}
%---------------------------------------------------------

Легко побачити, що остання вимога далеко не завжди може бути задовільнена. Дійсно, оскільки енергія замкнутої термодинамічної системи скінченна, то для деяких із фазових комірок значення $\vect{r}$ та $\vect{p}$ можуть виявитись такими, що значення енергії або декількох або навіть однієї молекули якщо вже вони опиняться у цій комірці, буде рівно значенню усієї енергії системи. Отже, усі інші молекули будуть змушені задовольнятися коміркою з нульовою енергією. Оскільки ця ситуація фатально малоймовірна, то, як ми побачимо далі, її здійсненням можна не перейматися. Виправданість цього базується на тих наслідках з методу Больцмана, що не раз підтверджувались на практиці. Крім того, існують й інші методи, такі як запропоновані Дарвіним та Фаулером, які обходять це ускладнення.

%% --------------------------------------------------------
\section{Наскільки малий може бути об’єм фазової комірки?}
%% --------------------------------------------------------

За сучасними уявленнями, які базуються на квантовій механіці, об’єм фазової комірки повинен бути кратним, а мінімальний в точності дорівнювати $h^3$. Це найменший фазовий об’єм, з яким оперує метод Больцмана. В принципі, розбиття на такі малі об’єми передбачає вже квантово-механічний опис поведінки системи, яке не припускає використовування понять координат і імпульсів одночасно. Але по суті хід міркувань у методі Больцмана не залежить від того яка точка зору прийнята – класична або квантова, хоча результати суттєво залежать від типу об’єктів (в даному випадку молекул).

%% --------------------------------------------------------
\section{Основне припущення прийняте статистики частинок}
%% --------------------------------------------------------

Об’єм, доступний молекулі у фазовому просторі визначається геометричним об’ємом, який займають молекули і значенням повної енергії системи $\varepsilon$. Це значення визначає границю області доступну для молекули у системі. Тоді пронумерувавши усі комірки, які потрапили у цю область та приписав їм деякі енергії $\varepsilon_i$, спробуємо розподілити молекули по цим коміркам так, щоб сума цих енергій дорівнювала $\varepsilon$. Це і є, за влучним зауваженням Шредінгера, основна задача статистики – розміщення об’єктів по припустимим станам системи при незмінних значеннях $\varepsilon$ і об’єму системи. Характер розміщення залежить від типу об’єктів. В статистиці Больцмана передбачається, що всі частинки (об’єкти) принципово різні, навіть якщо це і не так. Будемо говорити, що частинка знаходиться в $i$-тому стані, якщо вона попадає в $i$-ту комірку з енергією $\varepsilon_i$. Надалі будемо вважати, що в класичній границі в фазових комірках міститься дуже велика кількість комірок із мінімальним фазовим об’ємом $h^3$. Хаотичний рух частинок припускає абсолютно довільний розподіл частинок по коміркам так, що в $i$–й комірці знаходиться $n_i$ частинок. За Больцманом розподіл частинок по коміркам $n_1, n_2, n_3, \dots$ характеризує певний мікростан системи. Оскільки перестановки частинок всередині комірок не міняє повну енергію частинок і об’єм, який вони займають, то вона, очевидно, не міняє макростан системи, що характеризується макроскопічними параметрами, такими як об’єм, температура, тиск тощо. Міняючи розподіл частинок по коміркам можна отримати різні мікростани системи не міняючи її макростан.

Основне припущення прийняте у статистиці, що всі допустимі мікростани системи рівноімовірні. Справедливість цієї гіпотези неодноразово підтверджувалась виводами з неї отриманими, хай би яким інтуїтивно неймовірним може видаватися припущення, що мікростан, при якому всі частинки рівномірно розміщенні по всім допустимим станам і мікростан, при котрому одна частинка має енергію $\varepsilon$, а інші нічого, є рівноправні. В сутності, проте, нічого іншого a priori ми і не можемо припустити (Рис.~\ref{pic:Pic14.4}).

%---------------------------------------------------------
\begin{figure}[h!]\centering
	\includegraphics[width=0.5\linewidth]{Pic4}
	\caption{Два мікростани системи, що відповідають одному і тому ж макростану системи.}
	\label{pic:Pic14.4}
\end{figure}
%---------------------------------------------------------

%% --------------------------------------------------------
\section{Статистична вага мікростану}
%% --------------------------------------------------------

Нехай система містить $N$ частинок (наприклад молекул ідеального газу). При цьому зрозуміло:

\begin{equation}
	\sum n_i = N
	\label{eq:140.1}
\end{equation}

Кожний з можливих розподілів частинок по фазовим коміркам $\varepsilon_i$ характеризує певний макростан системи. Нехай деякий макростан характеризується $M$ комірками. Кількість різних розподілів з $N$ частинок по $M$ коміркам, що може реалізуватися в такому макростані системи дорівнює кількості мікростанів:

\begin{equation}
	W = \frac{N!}{n_1! \, n_2! \, \dots \, n_M!}
	\label{eq:140.2}
\end{equation}

Дійсно, кількість можливих комбінацій розміщення $N$ різних частинок так, щоб в першу комірку попало $n_1$ частинок, а в інші $N - n_1$ дорівнює:

\begin{equation}
	W_1 = \frac{N!}{n_1! \, (N - n_1)!} = C_N^{n_1}
	\label{eq:140.3}
\end{equation}

Кількість комбінацій розміщення $N - n_1$ частинок так, щоб у другій комірці опинилось $n_2$ частинок, а у інших $N - n_1 - n_2$ дорівнює кількості сполучень. Загальне число комбінацій таких, що в першій комірці знаходитьься $n_1$, у другій $n_2$, а у інших $N - n_1 - n_2$ дорівнює:

\begin{equation}
	W_{1,2} = \frac{N!}{n_1! \, (N - n_1)!} \cdot \frac{(N - n_1)!}{n_2! \, (N - n_1 - n_2)!} = \frac{N!}{n_1! n_2! (N - n_1 - n_2)!}
	\label{eq:140.4}
\end{equation}

Діючи так надалі по індукції отримаємо \eqref{eq:140.2}. Оскільки $N \sim 10^{23}$, то $W$ являє собою дуже велике число. Тому краще працювати з логарифмом $W$. Тоді:

\begin{equation}
	\ln W = \ln N! - \sum_i^M \ln (n_i)!
	\label{eq:140.5}
\end{equation}

При великих $n_i$ та $N$ доцільно користуватись формулою Стирлінга:

\begin{equation}
	N!={(2\pi N)}^{1/2}{\left(\frac{N}{e}\right)}^N
	\label{eq:140.6}
\end{equation}

Тоді

\begin{equation}
	\ln{\left(N!\right)}=\ \frac{1}{2}\ln{\left(2\pi N\right)}+NlnN-N\approx NlnN-N
	\label{eq:140.7}
\end{equation}

Підставляючи \eqref{eq:140.7} у \eqref{eq:140.5} отримуємо:

\begin{equation}
	\ln{W}=N\ln{N}-N-\sum_{i}^{M}{\left(n_i\ln{\left(n_i\right)}-n_i\right)=}\ N\ln{N}-\sum_{i}^{M}{n_i\ln{\left(n_i\right)}}.
	\label{eq:140.8}
\end{equation}

Де враховано рівняння \eqref{eq:140.1}. $W$ отримало назву \emph{статистичної ваги мікростану}.

%% --------------------------------------------------------
\section{Розподіл Максвела-Больцмана. Статистична сума.}
%% --------------------------------------------------------

Зрозуміло, що найбільш часто реалізований мікростан відповідає максимуму $W$, а відповідно і $\ln W$. Візьмемо варіацію \eqref{eq:140.8} з урахування умови:

\begin{equation}
	\sum_i^M \delta n_i = 0
	\label{eq:141.1}
\end{equation}

Тоді:

\begin{equation}
	\delta \ln W = - \sum_i^M \delta n_i \ln \frac{n_i}{N} - \sum_i^M \delta n_i = - \sum_i^M \delta n_i \ln \frac{n_i}{N}
	\label{eq:141.2}
\end{equation}

Вирази \eqref{eq:141.1} і \eqref{eq:141.2} сумісні тільки при рівномірному розподілу частинок по усім коміркам $n_1 = n_2 = \dots = n_M$ у повній відповідності до статистичної гіпотези. Проте, треба врахувати умову постійності енергії системи $U$:

\begin{equation}
	\sum{\varepsilon_i\delta n_i=\delta U=0},
	\label{eq:141.3}
\end{equation}

Щоб взяти варіацію $\ln W$ при одночасному виконуванні умов \eqref{eq:141.1} і \eqref{eq:141.3} скористаємося методом невизначених множників Лагранжа. Тоді вводячи $\alpha$ і $\beta$ (котрі треба буде визначити) варіюємо функцію Лагранжа $L\left(n_i\alpha,\beta\right)=\ln{W+\left(\alpha+1\right)\left(\sum n_i-N\right)+\beta\left(\sum{\varepsilon_in_i-U}\right)}$ отримаємо:

\begin{multline}\label{eq:141.4}
	-\sum\delta n_i(\ln{n_i}-1)-\left(\alpha+1\right)\sum\delta n_i-\beta\sum\varepsilon_i\delta n_i=\\=-\sum\delta n_i\left(\ln{n_i}+\alpha+\beta\varepsilon_i\right)=0.
\end{multline}
%
де $\alpha$ і $\beta$ треба обрати так, щоб для будь-якого $\delta n_i$ \eqref{eq:141.4} оберталося в нуль. Тоді:

\begin{equation}
	\ln{n_i}=-\alpha-\beta\varepsilon_i\ \ i\ \ n_i=e^{-\alpha}\cdot e^{-\beta\varepsilon_i}.
	\label{eq:141.5}
\end{equation}

Або

\begin{equation}
	n_i=n_0\cdot e^{-\beta\varepsilon_i}.
	\label{eq:141.6}
\end{equation}
%
де $n_0 = e^{-\alpha}$. Нам ще доведеться з’ясувати термодинамічний сенс $\alpha$ і $\beta$. Наразі зрозуміло, що розподіл \eqref{eq:141.6}, який отримав назву розподілу Максвела-Больцмана реалізує найбільше ймовірний мікростан системи. Підставляючи \eqref{eq:141.6} в \eqref{eq:140.8} знаходимо, що для $W_{\max}$:

\begin{equation}
	\ln{W_{\max}}=N\ln{N}+\alpha N+\beta U=\sigma\left(N,U\right)ю
	\label{eq:141.7}
\end{equation}

Тобто, максимальна кількість мікростанів залежить від параметрів, що характеризують макростан системи. Зауважимо, що самі числа заповнення $n_i$ так само залежать від $N$ і $U$.

Дійсно, в силу \eqref{eq:140.1}

\begin{equation}
	N=e^{-\alpha}\sum e^{-\beta\varepsilon_i}.
	\label{eq:141.8}
\end{equation}

Сума у \eqref{eq:141.8} називається сумою станів або статистичною сумою:

\begin{equation}
	Z_0=\sum e^{-\beta\varepsilon_i}.
	\label{eq:141.9}
\end{equation}

Згідно \eqref{eq:141.8} знаходимо, що $\alpha = \ln (Z_0 / N)$. Знаходимо, так само, що оскільки $Z_0 / N = e^{-\alpha}$, то:

\begin{equation}
	n_i=e^{-\alpha}\cdot e^{-\beta\varepsilon_i}=-\frac{N}{Z_0}\cdot\frac{\partial Z_0}{\partial\varepsilon_i}\cdot\frac{1}{\beta}=-\frac{N}{\beta}\cdot\frac{\partial\ln{Z_0}}{\partial\varepsilon_i}.
	\label{eq:141.10}
\end{equation}

\begin{equation}
	U = \sum_i^M \varepsilon_i n_i = \sum_i^M \varepsilon_i N \frac{e^{-\beta \varepsilon_i}}{Z_0} = N \frac{1}{Z_0} \sum_i^M \varepsilon_i e^{-\beta \varepsilon_i} = -N \frac{\partial \ln Z_0}{\partial \beta}
	\label{eq:141.11}
\end{equation}

Таким чином значення статистичної суми дозволяє знайти як числа заповнення, так і енергію системи. Забігаючи наперед зауважимо, що $Z_0$ дозволяє також знайти усі макроскопічні характеристики системи, тому основна задача статистики полягає в обрахування статистичної суми виходячи із певної моделі системи.

%% --------------------------------------------------------
\section{Максимальна статистична вага}\label{sec:MaxStatWeight}
%% --------------------------------------------------------

Для того, щоб з’ясувати сенс $\beta$, розглянемо дві системи які займають фіксовані об’єми в адіабатичній оболонці та знаходяться між собою у тепловому контакті. У рівновазі кожна з них характеризується значенням повної внутрішньої енергії $\hat{U}_1$ і $\hat{U}_2$. При цьому $\hat{U} = \hat{U}_1 + \hat{U}_2$. Проте, можливі й інші варіанти розподілу енергії між системами. Головне, щоб виконувалась умова $\hat{U} = U_1 + U_2$. Тоді повна кількість комбінацій припустимих станів складової системи буде

\begin{equation}
	W(N=N_1 + N_2, \hat{U})= \sum_{\hat{U}_1} W_1(N_1, U_1) W_2(N_2, \hat{U} - \hat{U}_1)
	\label{eq:142.1}
\end{equation}

Різноманітні припустимі конфігурації повної системи характеризується різними значеннями $U_1$. Найбільш ймовірна конфігурація має найбільшу кількість мікростанів і відповідає найбільшому члену у сумі \eqref{eq:142.1}:

\begin{equation}
	W_{\max} = W_1(N_1, \hat{U}_1) W_2(N_2, \hat{U} - \hat{U}_2)
	\label{eq:142.2}
\end{equation}

Виявляється, що для макроскопічних систем цей член приголомшливо більше ніж всі інші разом узяті.

Для того, щоб впевнитись у цьому, полегшимо собі задачу, зробивши оцінку для системи розділеної на дві частини з рівними об’ємами $V_1 = V_2 = V$ (Рис.~\ref{pic:Pic14.5}). Очевидно, що найбільш ймовірною конфігурацією буде така, коли $N_1 = N_2 = N/2$ і $\hat{U}_1 = \hat{U}_2 = \hat{U}/2$. Для оцінки простіше розглянути ситуацію коли енергія обох частин залишається постійною, а кількість частинок змінюється на $m$ (в одній половині збільшується, а у другій – зменшується). В загальному випадку $n$ частинок у об’ємі 1 можна розмістити $W_n$ способами:

\begin{equation}
	W_n = \frac{N!}{n! \, (N - n)!}
	\label{eq:142.3}
\end{equation}

Тоді імовірність спостерігати таку ситуацію складає:

\begin{equation}
	P_n = \frac{W_n}{\sum_n W_n} = \frac{W_n}{2^N} = \frac{N!}{n!\left(N-n\right)!}\left(\frac{1}{2}\right)^n\left(\frac{1}{2}\right)^{N-n}
	\label{eq:142.4}
\end{equation}

Де $(1/2)^n (1/2)^{N - n}$ є імовірність знайти $n$ частинок в об’мі 1 і $N - n$ в судині 2 після одного виміру. Зазначимо, що імовірність спостерігати хоч якусь конфігурацію становить:

\begin{equation}
	\sum_n P_n = \sum_{n=0}^N \left(\frac12\right)^2\left(\frac12\right)^{N - n} = 1
	\label{eq:142.5}
\end{equation}

Доречним буде питання: наскільки імовірним є відхилення від схеми розподілення частинок, яка трапляється найбільш часто, а саме:

\begin{equation}
	P_{N/2} = \frac{N!}{(N/2)! (N/2)!}\left(\frac12\right)^N
	\label{eq:142.6}
\end{equation}

Знайдемо імовірність такого відхилення $P_{N/2 - m}$ і віднесемо її до $P_{N/2}$:

\begin{equation}
	\frac{P_{N/2 - m}}{P_{N/2}} = \frac{(N/2)! \, (N/2)!}{(N/2 + m)! \, (N/2 - m)!}
	\label{eq:142.7}
\end{equation}

Логарифм відношення з урахуванням того, що зазвичай $m \ll N/2$ дасть:

%\begin{equation}
%	\ln \frac{P_{N/2 - m}}{P_{N/2}} \approx - \frac{2 m^2}{N}
%	\label{eq:142.8}
%\end{equation}

\begin{multline}\label{eq:142.8}
	\ln{\left(\frac{P_{\frac{N}{2}-m}}{P_\frac{N}{2}}\right)} = N\ln{\left(\frac{N}{2}\right)} - N - \left(\frac{N}{2}+m\right)\ln{\left(\frac{N}{2}+m\right)} + \left(\frac{N}{2}+m\right) - \\
	- \left(\frac{N}{2}-m\right)\ln{\left(\frac{N}{2}-m\right)} + \left(\frac{N}{2}-m\right) \cong \\
	N\ln{\left(\frac{N}{2}\right)} - \left(\frac{N}{2}+m\right)\left(\ln{\left(\frac{N}{2}\right)} + \ln{\left(1+\frac{2m}{N}\right)}\right) \\
	- \left(\frac{N}{2}-m\right)\left(\ln{\left(\frac{N}{2}\right)} + \ln{\left(1-\frac{2m}{N}\right)}\right) \cong \\
	\frac{2m}{N}\left(\frac{N}{2}-m - \frac{N}{2}-m\right) = -\frac{4m^2}{N}
\end{multline}

Це означає, що відношення ймовірностей дорівнює:

\begin{equation}\label{eq:142.9}
	e^{-\frac{4m^2}{N}}
\end{equation}

Нехай відносне відхилення числа частинок становить $2m/N = 10^{-10}$. Тоді власне $2m \approx 10^{23} \cdot 10^{-10} = 10^{13}$. Звідси відношення ймовірностей \eqref{eq:142.9} складе $\exp(-10^{26 + 23}) = e^{-1000} = 10^{-434}$. Це дуже мала величина. Якщо кількість частинок, які перетинають за одиницю часу умовну границю посередині судини оцінити як $Z = (1/4) n \langle v \rangle \approx 10^{23} \cdot 10^5 = 10^{28}$, тоді можна вважати, що в секунду газ $10^{28}$ разів змінює свою конфігурацію. Тоді якщо почекати $10^{434} \cdot 10^{-28} = 10^{406}$ секунд, то можна розраховувати на те, щоб побачити систему у такому стані. Оскільки вік Всесвіту порядку $10^{18}$ с, то це означає, що ми ніколи не зустрінемо такого стану. Зроблена оцінка виконана для кількості частинок але подібний результат вірний і для перерозподілу енергії між підсистемами і показує, що в сумі \eqref{eq:142.1} домінує найбільший член \eqref{eq:142.2} так, що \eqref{eq:142.1} фактично дорівнює \eqref{eq:142.2}.

%---------------------------------------------------------
\begin{figure}[h!]\centering
	\begin{subfigure}[b]{0.45\textwidth}
		\centering
		\includegraphics[width=\linewidth]{Pic5a}
		\caption{}
	\end{subfigure}
	\hfill
	\begin{subfigure}[b]{0.45\textwidth}
		\centering
		\includegraphics[width=\linewidth]{Pic5b}
		\caption{}
	\end{subfigure}
	\caption{Дві системи які займають фіксовані об’єми в адіабатичній оболонці та знаходяться між собою в а) тепловому та б) дифузійному контакті}
	\label{pic:Pic14.5}
\end{figure}
%---------------------------------------------------------

%% --------------------------------------------------------
\section{Формула Больцмана}
%% --------------------------------------------------------

Повернемося до розгляду повної системи. Для того, щоб знайти значення максимального члена в \eqref{eq:142.1} візьмемо варіацію \eqref{eq:142.2} по $U_1$, за умови, що $\delta U_1 = - \delta U_2$:

\begin{equation}
	\delta W=\left\{\left(\frac{\partial W_1}{\partial U_1}\right)_{N_1,V}W_2-\left(\frac{\partial W_2}{\partial U_2}\right)_{N_2,V\ }W_1\right\}\delta U_1=0
	\label{eq:143.1}
\end{equation}


\begin{wrapfigure}[18]{O}{0.2\linewidth}
	\centering
	\includegraphics[width=\linewidth]{Boltzmann}
	\captionsetup{justification=centering}
	\caption*{\href{https://en.wikipedia.org/wiki/Ludwig_Boltzmann}{Людвіг Больцман}\\(1844 -- 1906)}
\end{wrapfigure}
Звідси робимо висновок, що у рівновазі, що відповідає найбільш імовірній конфігурації системи:

\begin{equation}
	\left(\frac{\partial\ln{W_1}}{\partial U_1}\right)_{N_1,V}=\left(\frac{\partial\ln{W_2}}{\partial U_2}\right)_{N_2,V}
	\label{eq:143.2}
\end{equation}

З урахуванням явного виду виразу для $W_{\max}$ \eqref{eq:141.7} отримаємо, що:

\begin{equation}
	\beta=\left(\frac{\partial\ln{W}}{\partial U}\right)_{N,V}=\left(\frac{\partial\sigma\left(U,N\right)}{\partial U}\right)_{N,V}
	\label{eq:143.3}
\end{equation}

Зауважимо, що зробивши логічний перехід від \eqref{eq:143.2} до \eqref{eq:143.3}, ми неявно прийняли, що в тепловій рівновазі кожна з підсистем знаходиться у найбільш імовірній конфігурації, так само як і складова система. Але дві системи, які знаходяться у термодинамічній рівновазі мусять мати однакову температуру. Звідси $\beta$ має бути функцією температури. Як було показано методами феноменологічної термодинаміки для ізохорного процесу при $N = \text{const}$:

\begin{equation}
	dU = T \, dS = \frac{\theta}{k_\text{Б}}dS
	\label{eq:143.4}
\end{equation}
%
де $\theta$ --- енергетична температура. Порівнюючи \eqref{eq:143.4} і \eqref{eq:143.3}, приходимо до висновку, що $\beta = \theta^{-1} = (k T)^{-1}$ і $dS = k d \ln W = k d \ln W_{\max}$. Оскільки ентропія визначена с точністю до константи, то:

\begin{equation}
	S = k_\text{Б} \ln W
	\label{eq:143.5}
\end{equation}

Цей вираз отримав назву формули Больцмана. Вона розкриває фізичний сенс ентропії, як міри імовірності системи знаходитись у тому чи іншому стані.

%% --------------------------------------------------------
\section{Властивості ентропії за Больцманом}
%% --------------------------------------------------------

Визначення ентропії \eqref{eq:143.5} більш змістовно, ніж прийняте нами раніше оскільки дозволяє визначати у тому числі і ентропію нерівноважних станів системи розраховуючи $W$. При цьому другий закон термодинаміки втрачає свій абсолютний сенс і обертається у статистичний закон. Ентропія замкнутої системи може не тільки збільшуватись, але і зменшуватись. Проте, як показує зроблена нами оцінка, у переважній більшості випадків вона буде збільшуватись, навіть якщо система іноді і буде переходити у менш імовірний стан. Здоровий глузд підказує, що зрештою частіше будемо зустрічати стан системи, який може бути досягнуто більшою кількістю способів, а відтак і найбільш імовірний. Так, для двох систем приведених у контакт, якщо повна система переходить у більш імовірний стан, то:

\begin{equation}
	W_\text{кін} = W_1(\text{кін})W_2(\text{кін}) \geqslant W_1(\text{поч}) W_2(\text{поч})
	\label{eq:144.1}
\end{equation}

Знак рівності буде тільки тоді, якщо від початку системи були при однаковій температурі. Інтуїтивно прагнення ентропії до максимуму зрозуміло. Для кожної із систем накладаються обмеження на значення енергій $U_1$ і $U_2$. У складовій системі залишається тільки одне, накладене на повну енергію системи $U = U_1 + U_2$, що може тільки збільшити кількість дозволених станів. З \eqref{eq:144.1} слідує:

\begin{equation}
	S \geqslant S_1 + S_2
	\label{eq:144.2}
\end{equation}

Нескладно зрозуміти, що з \eqref{eq:143.5} і \eqref{eq:144.1} слідує і властивість адитивності ентропії. Зрозуміло, що властивість максимальності ентропії у стійкому рівноважному стані також слідує з визначення ентропії, як величини пропорційної імовірності знайти систему у тому чи іншому стані. Нарешті зазначимо, що таке визначення розкриває сенс третього закону термодинаміки згідно якому при $T \to 0$ ентропія прагне до скінченної границі, яка не залежить від параметрів, які характеризують стан системи (об’єм, тиск, поле, агрегатний стан і т.д.). Статистичне формулювання передбачає, що всі процеси при $T \to 0$ у системі зупиняються і вона переходить у єдино можливий припустимий стан. Якщо умовитися вважати ентропію цього стану рівною нулю, то усяка неоднозначність у визначенні цього стану зникає. Тоді визначення \eqref{eq:143.5} слід переписати як:

\begin{equation}
	S = k_\text{Б} \ln P
	\label{eq:144.3}
\end{equation}
%
де $P$ --- імовірність даного стану, що звісно ж пропорційно $W$ кількості припустимих станів.

%% --------------------------------------------------------
\section{Розподіл Максвела-Больцмана і числа заповнення}
%% --------------------------------------------------------

Розподіл Максвелла-Больцмана тепер може бути переписано як:

\begin{equation}
	n_i=e^{-\alpha}e^{-\frac{\varepsilon_i}{kT}}=const e^{-\frac{\varepsilon_i}{kT}}
	\label{eq:145.1}
\end{equation}

Значення константи може бути визначено із відношення нормування \eqref{eq:141.8}, що, зокрема, свідчить про її залежність від кількості частинок в системі. Проте більш продуктивно виразити її через якісь макроскопічні термодинамічні змінні, а не через статистичну суму. Зробимо це встановивши деякі важливі співвідношення. Для двох систем, які знаходяться у дифузному та тепловому контакті, максимальний член у \eqref{eq:142.1} знаходиться варіацією по $N_1$ і $U_1$:

\begin{multline}
	\delta W=\left\{\left(\frac{\partial W_1}{\partial U_1}\right)_{N_1,V}W_2-\left(\frac{\partial W_2}{\partial U_2}\right)_{N_2,V\ }W_1\right\}\delta U_1+ \\ +\left\{\left(\frac{\partial W_1}{\partial N_1}\right)_{U_1,V}W_2-\left(\frac{\partial W_2}{\partial N_1}\right)_{U_2,V\ }W_1\right\}\delta N_1=0
	\label{eq:145.2}
\end{multline}

Поділивши на $W_1 W_2$ отримаємо дві умови:

\begin{equation}
	{\frac{1}{W_1}\left(\frac{\partial W_1}{\partial U}\right)}_{N_1,V}=\frac{1}{W_2}\left(\frac{\partial W_2}{\partial U}\right)_{N_2,V\ }
	\label{eq:145.3}
\end{equation}

\begin{equation}
	\left(\frac{\partial\ln{W_1}}{\partial N}\right)_{U_1,V}=\left(\frac{\partial\ln{W_2}}{\partial N}\right)_{U_2,V}
	\label{eq:145.4}
\end{equation}

З \eqref{eq:145.3} слідує необхідність рівності температур в стані термодинамічної рівноваги. Друга ж умова еквівалентна тому, що в термодинамічній рівновазі дві системи будуть знаходитися у тепловому і дифузному контакті, якщо в них збігаються величини $(\partial \sigma / \partial N)_{V,U}$. Зауважимо, що цю величину можна зіставити з такою самою, яка визначається методами феноменологічної термодинаміки. Дійсно, за першим законом:

\begin{equation}
	T \, dS  = dU + p \, dV - \mu \, dN
	\label{eq:145.5}
\end{equation}

Звідки отримуємо, що:

\begin{equation}
	\frac{\mu}{T} = - \left( \frac{\partial S}{\partial N} \right)_{V,U}
	\label{eq:145.6}
\end{equation}

Зіставляючи з \eqref{eq:145.3} знайдемо, що:

\begin{equation}
	\left(\frac{\partial\ln{W}}{\partial N}\right)_{U,V}=-\frac{\mu}{\theta}
	\label{eq:145.7}
\end{equation}

Тепер зауважимо, що $\ln W_{\max}$ можна записати як:

\begin{equation}
	\ln{W_{\max}}=N\ln{N}+\left(\alpha-1\right)N+U/\theta
	\label{eq:145.8}
\end{equation}

Диференціюючи \eqref{eq:145.7} при $U = \const$ знайдемо:

\begin{equation}
	\alpha=-\frac{\mu}{\theta}-\ln{N}
	\label{eq:145.9}
\end{equation}

Звідси остаточно для розподілу Максвела-Больцмана знайдемо:

\begin{equation}
	n_i=N \exp\left\{\frac{\mu-\varepsilon_i}{\theta}\right\}
	\label{eq:145.10}
\end{equation}

З \eqref{eq:145.9} ми дістаємо висновку, що хімічний потенціал є як функцією температури, так і кількості частинок. Знак мінус в \eqref{eq:145.6} і \eqref{eq:145.7} вказує на те, що природнім ходом подій являється перехід частинок від системи с більшою кількістю частинок до меншої. Дійсно, для двох систем які знаходяться у тепловому та дифузному контакті:

\begin{equation}
	\delta S=\delta\left(S_1+S_2\right)=\left\{\left(\frac{{\partial S}_1}{\partial N_1}\right)_{U_1}-\left(\frac{\partial S}{\partial N_2}\right)_{U_2}\right\}\delta N=\left(-\frac{\mu_1}{T}+\frac{\mu_2}{T}\right)\delta N
	\label{eq:145.11}
\end{equation}

Отже, за умови $\mu_2 > \mu_1$ ентропія буде збільшуватися, якщо частинки переходять із системи $2$ у систему $1$, тобто $\delta N = \delta N_1 > 0$.

%% --------------------------------------------------------
\section{Фактори Гіббса і Больцмана. Велика статистична сума.}
%% --------------------------------------------------------

Знову повернемося до двох замкнутих систем, які знаходяться у дифузному та тепловому контакті. Але тепер одну із них будемо вважати термостатом, а другу власне системою (Рис.~\ref{pic:Pic14.6}). Повна енергія і кількість частинок термостата та системи залишаються незмінними $U_0$ та $N_0$. В рівновазі температура та хімічний потенціали їх збігаються. Зауважимо, що імовірність знайти систему із енергією $i$ та кількістю частинок $N_i$ пропорційна кількості можливих станів термостату. Дійсно, якщо фіксується стан системи, то кількість припустимих станів всього комплексу дорівнює кількості станів термостату. Отже імовірність $P(i, N_i)$ пропорційна числу припустимих станів термостату:

\begin{equation}
	P(\epsilon_i, N_i) \propto W_0(U_0 - \varepsilon_i, N_0 - N_i)
	\label{eq:146.1}
\end{equation}

Для відношень імовірностей двох довільних станів системи 1 та 2 маємо вираз:

\begin{equation}
	\frac{P_1(\epsilon_1, N_1)}{P_2(\epsilon_2, N_2)} = \frac{W_0(U_0 - \varepsilon_1, N_0 - N_1)}{W_0(U_0 - \varepsilon_2, N_0 - N_2)}
	\label{eq:146.2}
\end{equation}

Зауважимо тепер, що згідно \eqref{eq:141.7}:

\begin{equation}
	W\left(U_0,N_0\right)=W_{\max}\equiv\exp{\sigma}\left(U_0,N_0\right)
	\label{eq:146.3}
\end{equation}

Тому відношення \eqref{eq:146.2} можна переписати як:

\begin{equation}
	\frac{P\left(\varepsilon_1,N_1\right)}{P\left(\varepsilon_2,N_2\right)}=\frac{\exp{\sigma\left(U_0-\varepsilon_1,N_0-N_1\right)}}{\exp{\sigma\left(U_0-\varepsilon_2,N_0-N_2\right)}}=\exp\Delta\sigma
	\label{eq:146.4}
\end{equation}
%
де різниця безрозмірних ентропій (виражених із точністю до сталої Больцмана) може бути розвинута в ряд поблизу свого неймовірно гострого максимуму при $U_0$ та $N_0$ по $\epsilon_1-\epsilon_2$ та $N_1 - N_2$:

\begin{equation}
	\Delta\sigma = -\left(\frac{\partial\sigma}{\partial\epsilon}\right)_{N_0} (\epsilon_1 - \epsilon_2) -
	\left(\frac{\partial\sigma}{\partial N}\right)_{U_0} (N_1 - N_2) = \frac{\mu(N_1 - N_2)}{\theta}
	- \frac{\epsilon_1 - \epsilon_2}{\theta}
	\label{eq:146.5}
\end{equation}

Відношення імовірностей є відношенням експоненційних членів, що отримали назву фактору Гібса:

\begin{equation}
	\exp\left(\frac{\mu N}{\theta}-\frac{\varepsilon}{\theta}\right)
	\label{eq:146.6}
\end{equation}

Для системи із сталим числом частинок перший член в показнику експоненти випадає і відношення імовірностей визначається \emph{фактором Больцмана}.

Для того, щоб перейти від відносних імовірностей до абсолютних слід знайти нормуючий множник $C$, що визначається із умови нормування: імовірність знайти систему в хоч якомусь із станів має дорівнювати 1, оскільки реалізація кожного із станів є взаємовиключні події:

\begin{equation}
	C\sum_{i}\exp\left(\frac{\mu N_i}{\theta}-\frac{\varepsilon_i\left(N\right)}{\theta}\right)=1
	\label{eq:146.7}
\end{equation}

Сума в \eqref{eq:146.7} називається великою статистичною сумою $Z$. Із урахуванням її визначення абсолютні значення імовірностей становлять:

\begin{equation}
	P\left(\varepsilon_i,N\right)=\frac1Z \exp\left(\frac{\mu N_i}{\theta}-\frac{\varepsilon_i\left(N\right)}{\theta}\right).
	\label{eq:146.8}
\end{equation}

Отриманий вираз дає можливість отримати середні значення термодинамічних величин. Зокрема середня енергія системи (це власне є внутрішня енергія $U$, з якою ми маємо справу в термодинаміці), що знаходиться у дифузному та тепловому контакті із термостатом становить:

\begin{equation}
	U =\sum_{i}{\varepsilon_iP\left(\varepsilon_i,N\right)}=\sum_{i}{\varepsilon_i\frac{\exp\left(\frac{\mu N_i}{\theta}-\frac{\varepsilon_i\left(N\right)}{\theta}\right)}{Z}}=\frac{\theta^2}{Z}\frac{\partial Z}{\partial\theta}=\theta^2\frac{\partial\ln{Z}}{\partial\theta}.
	\label{eq:146.9}
\end{equation}

Отриманий вираз збігається із \eqref{eq:141.10}, якщо взяти до уваги дві обставини. Перша, очевидна: це те, що $\beta = -1$. Друга слідує із того, що при незмінній кількості частинок $N$ велика статистична сума $Z$ перетворюється на статистичну суму $Z_0$. Формально, це видно із того, що в такому разі сума складається із факторів Больцмана, а отже підсумовуються фактори, пропорційні числам заповнення мікростану. Більш фізичне пояснення полягає в тому, що розподіл частинок по енергіям не залежить від кількості частинок, а визначається тільки температурою і відповідна середня енергія визначає середнє значення для окремої частинки, яке аж ніяк не може залежати від кількості частинок. Тому останнє слід множити на кількість частинок $N$, щоб визначити повну енергію системи $U$.

Натомість вираз \eqref{eq:146.8} дає можливість визначити середнє значення частинок в системі, що знаходиться у дифузному та тепловому контакті із термостатом:

\begin{equation}
	\langle N \rangle = \sum_i N_i P(\epsilon_i, N) = \sum_i N_i \sum_{i}{N_i\frac{\exp\left(\frac{\mu N_i}{\theta}-\frac{\varepsilon_i\left(N\right)}{\theta}\right)}{Z}} = \frac{\theta}{Z} \frac{\partial Z}{\partial\mu}.
	\label{eq:146.10}
\end{equation}

%---------------------------------------------------------
\begin{figure}[h!]\centering
	\includegraphics[width=0.5\linewidth]{Pic6}
	\caption{До ілюстрації понять фактор Гібса і фактор Больцмана}
	\label{pic:Pic14.6}
\end{figure}
%---------------------------------------------------------

%% --------------------------------------------------------
\section{Формула Больцмана. Експеримент Перрена}
%% --------------------------------------------------------

При виведенні розподілу Максвелла-Больцмана ми ніяк не конкретизували вираз для енергії окремої частинки. У загальному випадку її можна подати у вигляді:

\begin{equation}
	\varepsilon_i = \varepsilon' + \varphi(\vect{r}).
	\label{eq:147.1}
\end{equation}
%
де $\varphi$ --- потенціал сил діючих на молекули, а $\varepsilon$ частина повної енергії незалежної від положення молекули. Вона включає кінетичну енергію, обертання, коливання, збудження атома і т.п. Тоді число станів $n_i$ згідно \eqref{eq:145.1} може бути подано як:

\begin{equation}
	n_i=e^\frac{-\alpha-\varphi\left(\vec{r}\right)}{\theta}e^{-\frac{{\varepsilon'}_i}{kT}}
	\label{eq:147.2}
\end{equation}

Підсумовуючи по частині фазового об’єму (практично може бути зручнішим інтегрування), для числа частинок в елементі об’єму отримуємо:

\begin{equation}
	n^\prime=e^\frac{-\alpha-\varphi\left(\vec{r}\right)}{\theta}\sum_{i} e^{-\frac{{\varepsilon'}_i}{\theta}}
	\label{eq:147.3}
\end{equation}

\begin{wrapfigure}[16]{O}{0.2\linewidth}
	\centering
	\includegraphics[width=\linewidth]{Perrin}
	\captionsetup{justification=centering}
	\caption*{\href{https://en.wikipedia.org/wiki/Jean_Baptiste_Perrin}{Жан Батіст Перрен}\\(1870 -- 1942)}
\end{wrapfigure}
Домноживши на масу молекули $m$, отримуємо вираз еквівалентний \eqref{eq:147.3}, але для густини речовини. Якщо позначити за $\rho_0$ густину молекул при $\varphi = 0$, тобто величину $\rho_0 = e^{-(\alpha + 1)/\theta} \Sigma_i \exp(-\varepsilon_i / \theta)$, то остаточно отримаємо формулу Больцмана для густини розподілу густини газу при постійній температурі:

\begin{equation}
	\rho(\vect{r})=\rho_0e^{-\frac{\varphi(\vect{r})}{kT}}
	\label{eq:147.4}
\end{equation}


Формула дозволяє знайти розподіл частинок в будь-якому потенціальному полі. Використовуючи цю формулу Жан Батіст Перрен визначив експериментально значення постійної Больцмана досліджуючи розподіл суспензії частинок гуммігута у воді. Якщо розглядати ці частинки, як газ макромолекул, то поле сили тяжіння діюче на таку макромолекулу складає $(m - m_{\text{ж}}) g z$, де $z$ --- вісь співнапрямлена з $g$. Тоді, згідно \eqref{eq:147.4}:

\begin{equation}
	\rho(z) = \rho_0 e^{-(m - m_{\text{р}}) g z / k_\text{<} T}
	\label{eq:147.5}
\end{equation}

Вимірюючи $\rho(z)$ можа знайти $k_{\text{Б}}$.

%---------------------------------------------------------
\begin{figure}[h!]\centering
	\includegraphics[width=0.5\linewidth]{Pic7}
	\caption{До ілюстрації розподілу суспензії частинок гуммігута у воді в експерименті Жана Батіста Перрена}
	\label{pic:Pic14.7}
\end{figure}
%---------------------------------------------------------

%% --------------------------------------------------------
\section{Барометрична формула}
%% --------------------------------------------------------

В окремому випадку дія гравітаційного поля Землі, розподіл Больцмана може бути отриманий з умови механічної рівноваги стовпця газу у полі сили тяжіння (Рис.~\ref{pic:Pic14.8}).

Маса стовпця газу, який спирається на площинку одиничної площини $S$ і заключений між $h$ і $h + dh$ складає $n m dh$. Вага цього стовпця повинна бути врівноважена різницею тисків $dP = (dP / dh) dh$. Звідси отримуємо:

\begin{equation}
	\frac{dP}{dh} = - n m g
	\label{eq:148.1}
\end{equation}

Але, згідно основній формулі молекулярно-кінетичної теорії газів $P = n k T$, звідки \eqref{eq:148.1} переписується як:

\begin{equation}
	kT\frac{dn}{dh}=-nmg.
	\label{eq:148.2}
\end{equation}

Інтегруючи отримуємо:

\begin{equation}
	n=\const\cdot e^{-\frac{mgh}{RT}}=\const\cdot e^{-\frac{\mu gh}{RT}}
	\label{eq:148.3}
\end{equation}

Константа в загальному випадку визначаються з умови нормування. Остаточно для газу у полі тяжіння:

\begin{equation}
	n = n_0 e^{-m g h / k T}
	\label{eq:148.4}
\end{equation}
%
де $n_0$ значення $n$ при $h = 0$. Переходячи від $n$ до $P$ отримаємо для тиску:

\begin{equation}
	P = P_0 e^{-\mu g h / R T}
	\label{eq:148.5}
\end{equation}

Це тиск газу, як функція висоти --- барометрична формула. Слід зауважити, що такий вивід формули Больцмана може призвести до помилкових результатів для суміші газів. Якщо підійти до \eqref{eq:148.5} формально, то для суміші під $\mu$ можна необґрунтовано взяти середню молекулярну вагу суміші. Проте, оскільки поведінка різних газів незалежна, то під $P$ в \eqref{eq:148.5} слід мати на увазі парціальний тиск газу. Звідси слідує, що концентрація більш важких газів зменшується із висотою швидше, ніж легких. На цьому принципи базується метод седиментації метод осадження більш важких фракцій, наприклад в суспензіях порошків.

%---------------------------------------------------------
\begin{figure}[h!]\centering
	\includegraphics[width=0.5\linewidth]{Pic8}
	\caption{До виводу барометричної формули}
	\label{pic:Pic14.8}
\end{figure}
%---------------------------------------------------------

%% --------------------------------------------------------
\section{Вивід розподілу Максвела із розподілу Максвела-Больцмана}
%% --------------------------------------------------------

Чи можна отримати розподіл Максвела із розподілу Максвела-Больцмана? Для того, щоб це зробити перепишемо \eqref{eq:145.1} як:

\begin{equation}
	n_i=e^{-\alpha}e^{-\left(p_{ix}^2+p_{iy}^2+p_{iz}^2\right)/2mkT}
	\label{eq:149.1}
\end{equation}
%
та підсумувавши по просторовій частині фазового об’єму, отримаємо кількість частинок в елементі фазового об’єму $\Delta p_{ix} \Delta p_{iy} \Delta p_{iz}$:

\begin{multline}
	n^{\prime\left({\vect{p}}_i\right)}=\frac{V}{\Delta x \Delta y \Delta z}e^{-\alpha}\exp\left(-\frac{ p^2_{ix} + p^2_{iy} + p^2_{iz }}{2mkT}\right) = \\ = \const\cdot\exp\left(-\frac{ p^2_{ix} + p^2_{iy} + p^2_{iz }}{2mkT}\right).
	\label{eq:149.2}
\end{multline}
%
де $n^{\prime\left({\vect{p}}_i\right)}$ --- кількість частинок з енергією $\varepsilon_i$ у елементі $\Delta p_{ix} \Delta p_{iy} \Delta p_{iz}$.

Тепер, згідно уявленням, які ми використали при виведенні розподілу Максвела, нас буде цікавити густина розподілу $F(\vect{p}_i)$, що визначається через імовірність знайти в елементі $\Delta p_{ix} \Delta p_{iy} \Delta p_{iz}$ певну кількість частинок як:

\begin{equation}
	\frac{n^{\prime\left({\vec{p}}_i\right)}}{N} \Delta p_{ix} \Delta p_{iy} \Delta p_{iz} = F(\vect{p}_i) \Delta p_{ix} \Delta p_{iy} \Delta p_{iz}.
	\label{eq:149.3}
\end{equation}
%
де $N$ --- повна кількість частинок. Співставляючи \eqref{eq:149.2} і \eqref{eq:149.3} бачимо, що:

\begin{equation}
	F\left({\vect{p}}_i\right)=C\exp{\left(-\frac{p_x^2+p_y^2+p_z^2}{2mkT}\right)}
	\label{eq:149.4}
\end{equation}
%
де константа $C$ може бути отримана підсумовуванням \eqref{eq:149.4} по імпульсній частині фазового простору. Перейшовши від підсумовування до інтегрування щільності розподілу, отримаємо одиницю $\int{F\left(\vect{p}\right)d^3\vec{p}=1}$. Звідси остаточно:

\begin{equation}
	F(\vect{p}) = \frac{\exp{\left(-\frac{p_x^2+p_y^2+p_z^2}{2mkT}\right)}}%
	{
		\int\limits_{-\infty}^{\infty} e^{-\frac{p_x^2}{2mkT} }dp_x
		\int\limits_{-\infty}^{\infty} e^{-\frac{p_y^2}{2mkT} }dp_y
		\int\limits_{-\infty}^{\infty} e^{-\frac{p_z^2}{2mkT} }dp_z
	}
	\label{eq:149.5}
\end{equation}

Інтегрування дасть у знаменнику $(2 \pi m k T)^{3/2}$. Тоді:

\begin{equation}
	F\left(\vect{p}\right)=(2 \pi m k T)^{3/2} \exp\left(-\frac{p_x^2+p_y^2+p_z^2}{2mkT}\right)
	\label{eq:149.6}
\end{equation}
%
що і є розподілом Максвела с тою формальною різницею, що для імпульсів $\vect{p}$. Саме тому в експоненті $m$ стоїть у чисельнику, а в константі перед експонентою в знаменнику.

%% --------------------------------------------------------
\section{Флуктуації термодинамічних величин}
%% --------------------------------------------------------

Зупинимося детальніше на питанні флуктуацій термодинамічних величин. Згідно основній ідеї наведеного розгляду будь-яка фізична характеристика, що фіксується експериментально, є тільки середнім значенням, що відповідає найбільш імовірній конфігурації системи. Правомірно задатися питанням наскільки великі можуть бути флуктуації величин і чи не є вони настільки великими, що статистичний розгляд втрачає будь-який сенс.

Мірою флуктуації прийнято називати корінь із середньоквадратичного відхилення:

\begin{equation}
	\langle \delta f\rangle = \sqrt{\langle \Delta f^2 \rangle} = \sqrt{ \langle (f  - \langle f \rangle)^2}
	\label{eq:150.1}
\end{equation}

Розкриваючи вираз під коренем праворуч отримаємо

\begin{equation}
	\langle (f  - \langle f \rangle)^2 = \langle f^2 \rangle - 2 \langle f \rangle\langle f \rangle + \langle f \rangle^2
	\label{eq:150.2}
\end{equation}

Таким чином для знаходження флуктуації деякої фізичної величини $f$ слід знати її середнє значення і середній квадрат.

Так середня енергія молекули газу, згідно розподілу Максвела-Больцмана, може бути знайдена як:

\begin{equation}
	\langle \varepsilon \rangle = \sum_i \frac{n_i}{N} \varepsilon_i
	= \frac{\sum_i \varepsilon_i e^{-\alpha} e^{-\beta \varepsilon_i}}{\sum e^{-\alpha} e^{-\beta \varepsilon_i}}
	= -\frac{\partial \ln Z_0}{\partial \beta}
	\label{eq:150.3}
\end{equation}

Аналогічно для середнього квадрату:

\begin{equation}
	\langle \varepsilon^2 \rangle = \frac{\sum_i \varepsilon_i^2 e^{-\alpha} e^{-\beta \varepsilon_i}}{\sum e^{-\alpha} e^{-\beta \varepsilon_i}}
	= \frac{1}{Z_0} \frac{\partial^2 Z_0}{\partial \beta^2}
	\label{eq:150.4}
\end{equation}

Тоді середньоквадратична флуктуація визначається як:

\begin{equation}
	\langle \Delta \varepsilon^2 \rangle = \frac{1}{Z_0} \frac{\partial^2 Z_0}{\partial \beta^2} - \frac{1}{Z_0^2} \left( \frac{\partial Z_0}{\partial \beta} \right)^2
	= \frac{\partial^2 \ln Z_0}{\partial \beta^2}
	= -\frac{\partial \langle \varepsilon \rangle}{\partial \beta}
	\label{eq:150.5}
\end{equation}

Для $N$ молекул газу ця величина повинна бути помножена на $N$. Тоді враховуючи, що $U = N \langle \varepsilon \rangle$, а $\beta = (k_\text{Б} T)^{-1}$ маємо:

\begin{equation}
	\langle \Delta E^2 \rangle = \langle \Delta U^2 \rangle = - \frac{\partial \langle U \rangle}{\partial \beta} = kT^2 \frac{\partial U}{\partial T} = C_V kT^2
	\label{eq:150.6}
\end{equation}

Тобто квадратичну флуктуацію енергії можна встановити із макроскопічних термодинамічних величин. Так, для одного моля одноатомного газу $U = (3/2) N k T$ і флуктуація становить:

\begin{equation}
	\frac{\langle \delta U \rangle}{\langle U \rangle} = \sqrt{\frac{2}{3N}} \approx 10^{-12}
	\label{eq:150.7}
\end{equation}

Звісно це незначна величина. Проте для нанорозмірних об’єктів флуктуація може виявитися значною. Знаючи флуктуації енергії \eqref{eq:150.6}, можна виразити флуктуації $P, V, S, T$ та інших термодинамічних змінних.

Для того, щоб знайти флуктуації чисел заповнення $n_i$ спочатку отримаємо вираз для флуктуації кількості молекул $N$. Для цього скористаємося виразом для великої статистичної суми. Звідки знаходимо:

\begin{equation}
	\langle N \rangle = \sum_i N_i P(\varepsilon_i, N) = \sum_{N,i} N_i \frac{\exp\left( \frac{\mu N_i}{\theta} - \frac{\varepsilon_i (N)}{\theta} \right)}{Z} = \frac{\theta \, \partial Z}{Z \, \partial \mu} = \theta \frac{\partial \ln Z}{\partial \mu}
	\label{eq:150.8}
\end{equation}

\begin{equation}
	\langle N^2 \rangle = \sum_{N_i} N_i^2 \frac{\exp\left( \frac{\mu N}{\theta} - \frac{\varepsilon_i (N)}{\theta} \right)}{Z} = \frac{\theta^2 \, \partial^2 Z}{Z \, \partial \mu^2}
	\label{eq:150.9}
\end{equation}

Діючи так само як і при виведенні \eqref{eq:150.5} отримаємо:

\begin{equation}
	\langle \Delta N^2 \rangle = \frac{\theta^2 \, \partial^2 Z}{Z \, \partial \mu^2} - \left( \frac{\theta \, \partial Z}{Z \, \partial \mu} \right)^2 = \theta \frac{\partial}{\partial \mu} \left( \frac{\theta \, \partial Z}{Z \, \partial \mu} \right) = \theta \frac{\partial \langle N \rangle}{\partial \mu}
	\label{eq:150.10}
\end{equation}

Тепер скористаємося тим, що числа заповнення $n_i$ самі по собі є середніми значеннями. Тому для оцінки їх флуктуації може бути застосований останній вираз:

\begin{equation}
	\langle \Delta n_i^2 \rangle = \theta \frac{\partial n_i}{\partial \mu}
	\label{eq:150.11}
\end{equation}

Диференціюючи вираз \eqref{eq:145.9} по $\mu$ знайдемо:

\begin{equation}
	\langle \Delta n_i^2 \rangle = n_i
	\label{eq:150.12}
\end{equation}

Відносна флуктуації чисел заповнення тоді становить:

\begin{equation}
	\frac{\langle \delta n_i \rangle}{\langle n_i \rangle} = \frac{1}{\sqrt{n_i}}
	\label{eq:150.13}
\end{equation}

З огляду на пропорційність чисел заповнення повній кількості частинок можемо очікувати, що відносна флуктуація кількості частинок теж зворотньо пропорційна кореню із $N$:

\begin{equation}
	\frac{\langle \delta N \rangle}{\langle N \rangle} = \frac{1}{\sqrt{N}}
	\label{eq:150.14}
\end{equation}

Вирази \eqref{eq:150.7}, \eqref{eq:150.13} і \eqref{eq:150.14} є наслідком закону великих чисел і добре знайомі кожному, хто займався статистикою. Звідси видно, що для макроскопічних систем наслідки термодинаміки і статистики добре обґрунтовані. Інший спосіб отримання формули \eqref{eq:150.14} наведено в Додатку \ref{add:B}.

\Closesolutionfile{answer}